
\section{Related Works}
\label{sec:related}

\noindent
\textbf{Acoustic Side-Channel Attacks (ASCAs)} exploit sound emanations from electronic devices, particularly keyboards, to infer sensitive information like passwords and PINs. With the ubiquity of microphones, these attacks have evolved over decades, addressing prior limitations while presenting new challenges and research opportunities. This section reviews their progression, key challenges, and areas for improvement.

% Traditional methods: 7 citations
\textbf{Traditional statistical ASCA approaches} relied on signal processing and acoustic feature analysis, using methods like cosine similarity and cross-correlation. Cosine similarity effectively modeled inter-keystroke timings to infer numeric inputs such as PIN codes~\cite{liu2019human}, while cross-correlation highlighted acoustic signal consistency across similar keystrokes~\cite{panda2020behavioral}. Techniques such as Time Difference of Arrival (TDoA) improved accuracy by leveraging geometric keyboard positions~\cite{de2019differential, zhu2014context}. Hidden Markov Models (HMMs), enhanced by timing and language-based post-processing (e.g., n-grams and spelling correction), further reconstructed text from keystroke acoustics~\cite{foo2010timing, zhuang2009keyboard}. Nevertheless, these methods lacked robustness, performed well primarily on limited numeric or predictable text, and struggled significantly with full alphanumeric recovery under realistic noise and typing variations~\cite{taheritajar2024survey}.
% \textbf{Traditional statistical approaches in ASCA research} relied on signal processing and acoustic feature analysis, utilizing methods such as cosine similarity and cross-correlation to characterize keystrokes. Cosine similarity-based approaches showed potential in modeling inter-keystroke timings to infer short numeric inputs such as PIN codes ~\cite{liu2019human}, while cross-correlation methods provided insights into acoustic signals' consistency across similar keystrokes ~\cite{panda2020behavioral}. Localization of keystrokes based on methods like Time Difference of Arrival (TDoA) further enhanced accuracy by exploiting geometric positional differences on keyboards ~\cite{de2019differential, zhu2014context}. Moreover, Hidden Markov Models (HMMs) combined timing analysis and language-based post-processing (e.g., n-gram language models and spelling correction) to reconstruct coherent text from keystroke acoustics ~\cite{foo2010timing, zhuang2009keyboard}. However, these statistical approaches lacked sufficient robustness, typically performing well only on limited numeric or predictable texts. They struggled significantly with full-alphanumeric keyboard recovery and effectively degraded under even minor variations in real-world noise and typing behaviors ~\cite{taheritajar2024survey}.

% Machine learning methods: 7 citations
\textbf{Machine learning methods} enhanced ASCA robustness by using advanced models like clustering, Support Vector Machines (SVMs), and Hidden Markov Models (HMMs). K-means clustering combined with Mel-Frequency Cepstrum Coefficients (MFCCs) effectively grouped keystrokes into distinct acoustic classes~\cite{wit2014all}. Systematic use of HMMs~\cite{zhuang2009keyboard}, SVMs~\cite{wang2016accurate}, logistic regression, and random forests~\cite{anand2018keyboard} further improved accuracy in controlled environments. Despite these advances, machine learning methods remained highly sensitive to typing styles, keyboard variations, and acoustic noise, with accuracy dropping by 30–50\% under realistic conditions~\cite{halevi2012closer, fang2018no, liu2015snooping}.
%\textbf{Machine learning methods} improved robustness by leveraging advanced data-driven models, with clustering techniques, Support Vector Machines (SVMs), and Hidden Markov Models (HMMs) enhancing keystroke recognition accuracy. Clustering methods such as K-means, combined with Mel-Frequency Cepstrum Coefficients (MFCCs), successfully grouped and classified keystrokes into distinct acoustic classes, improving recognition rates considerably ~\cite{wit2014all}. Similarly, more systematic use of HMMs ~\cite{zhuang2009keyboard}, SVMs~\cite{wang2016accurate}, logistic regression, and random forest models~\cite{anand2018keyboard} offered promising performance under controlled indoor environments. These solutions took advantage of more informative acoustic features and richer statistical modeling, achieving marked improvements in keystroke recognition accuracy. Nevertheless, machine learning approaches remained highly sensitive to variations in typing styles, keyboard types, and especially acoustic noise and room conditions. Consequently, their accuracy frequently deteriorated as environmental realism increased—often dropping significantly, with reductions of 30–50\% in realistic noisy scenarios ~\cite{halevi2012closer, fang2018no, liu2015snooping}.

% Deep Learning: 6 citations
\textbf{Recent deep learning (DL) architectures} significantly outperformed traditional methods by extracting hierarchical features from acoustic spectrograms. CNNs, ConvMixer models, and CNN-RNN hybrids achieved over 90\% keystroke classification accuracy under clean conditions~\cite{harrison2023practical, akinbi2023password, giallanza2019keyboard, toreini2015acoustic}; for instance, Harrison et al. reached above 93\% accuracy using CoAtNet models~\cite{harrison2023practical}. However, these DL methods remain constrained by their need for extensive labeled data, high computational costs, and notably poor performance (often below 40\%) in realistic noisy environments~\cite{giallanza2019keyboard}. This fundamental lack of robustness arises from purely local classification approaches, failing to leverage global contextual semantics, thus limiting their practical effectiveness in real-world ASCA scenarios.
%\textbf{More recently, Deep Learning (DL) architectures} significantly surpassed traditional methods due to their superior capability for hierarchical feature extraction and meaningful pattern representation from acoustic spectrograms. Notably, Convolutional Neural Networks (CNNs), ConvMixer models, and CNN-RNN hybrids have demonstrated state-of-the-art performance, achieving keystroke classification accuracy surpassing 90\% in clean dataset conditions ~\cite{harrison2023practical, akinbi2023password, giallanza2019keyboard, toreini2015acoustic}. For instance, Harrison et al. achieved over 93\% classification accuracy on clean data using architectures combining convolution with self-attention (CoAtNet) ~\cite{harrison2023practical}. Despite these significant accuracy gains, deep learning methods still suffer from fundamental issues: they require large volumes of labeled training data, substantial computational resources for training, and their performance sharply declines (often below 40\% accuracy) when confronted with realistic noisy environments ~\cite{giallanza2019keyboard}. Thus, even state-of-the-art DL architectures struggle to reliably recover sensitive text in practical side-channel scenarios, particularly passwords, where small classification errors render entire information recovery attempts futile. Therefore, the fundamental unresolved problem is the lack of robustness to acoustic noise, a consequence of purely local classification decisions without context. First-generation statistical methods performed only local comparisons (e.g., acoustic similarity between isolated keystrokes), while second-generation ML and DL methods, although considering richer representations, still predominantly performed local classification, failing to leverage global correlations among keys or contextual semantics. This locality severely limited all previous approaches in noisy real-world settings.

\textbf{Transformer architectures} recently emerged as powerful methods capable of modeling global relationships and handling long-range dependencies, significantly improving noise robustness. Although transformers have demonstrated strong results in computer vision (Vision Transformers, VTs) and natural language processing (Large Language Models, LLMs), their application to acoustic side-channel attacks remains unexplored. VTs, through self-attention mechanisms, can identify global acoustic patterns, potentially addressing local noise challenges, while LLMs' proven contextual reasoning abilities~\cite{dosovitskiy2020image, brown2020language} could correct misclassifications in noisy keystroke scenarios. Nevertheless, transformer implementation in ASCA faces practical challenges, primarily computational costs, though recent lightweight fine-tuning methods like LoRA and QLoRA significantly alleviate these concerns~\cite{hu2021lora, dettmers2024qlora}. Yet, the effectiveness of these lightweight methods specifically for acoustic keystroke error correction has not been validated.

%\textbf{Meanwhile, transformer architectures} have emerged as a powerful new class of solutions strongly tailored to global relationship modeling, capable of effectively handling long-range dependencies and noise robustness. While extensively validated in computer vision (as Vision Transformers, VTs) and natural language domains (as Large Language Models, LLMs), transformers remain unexplored in acoustic keyboard side-channel scenarios. VTs leverage a self-attention mechanism capable of identifying salient global acoustic features across entire spectrogram input, and thus offer excellent potential to overcome local noise issues. Concurrently, LLMs have exceptional contextual reasoning capabilities, proven in numerous language-related correction tasks, which could significantly strengthen ASCA pipelines by correcting misclassifications or typos when classifying text from noisy keystroke acoustics ~\cite{dosovitskiy2020image, brown2020language}. Despite this strong potential, no prior ASCA research has tested whether transformers' global correlation capability can solve the local degradation problem in acoustic side-channel classification or leveraged LLMs for contextual error correction under serious acoustic interference.

%\textbf{Practical challenges in implementing transformers for ASCA} include computational costs, but emerging lightweight fine-tuning techniques like LoRA and QLoRA could provide feasible solutions. Recently emerged lightweight fine-tuning methods, specifically LoRA and Quantized LoRA (QLoRA), offer tangible solutions by drastically reducing computational costs associated with fine-tuning and inference of LLMs on a specific domain ~\cite{hu2021lora, dettmers2024qlora}. Though these techniques provide plausible methods for feasible ASCA deployment, their performance and efficiency specifically in acoustic keystroke error-correction have not been validated.

\textbf{In summary}, current ASCA methods are limited by their reliance on local acoustic classification, greatly reducing accuracy in noisy environments. Transformers, capable of modeling global dependencies and context, represent a promising yet unexplored solution to these limitations. Recent lightweight fine-tuning advancements (LoRA, QLoRA) further suggest practical feasibility. Here, we explicitly evaluate the hypothesis that integrating vision transformers and LLM-based contextual correction addresses noise robustness without sacrificing computational efficiency, conclusively determining their value in practical ASCA scenarios.

%\textbf{In summary,} existing ASCA approaches face challenging limitations arising primarily from their reliance on local acoustic classification, severely impacting accuracy under realistic, noisy environments. Transformers, with their unique ability to model global dependencies and contextual relationships, represent an untested yet highly promising solution able to effectively address these inherent limitations. Moreover, recent advances in lightweight LLM fine-tuning (LoRA, QLoRA) suggest that this powerful solution could be practically viable. In this work, we explicitly test and verify the hypothesis that integrating vision transformers and LLM-based global contextual correction can overcome existing limitations in ASCA—specifically the robustness to real-world noise—while minimally impacting computational feasibility. Our results will conclusively demonstrate whether transformer-enhanced pipelines represent the much-needed advancement to robust acoustic side-channel attacks on keyboards.
